\section{Finite-Fock Expectations and Comparison Limits}
\label{sec:ch19-finite-comparison}

\subsection{The comparison problem: the same formula is not the same state}

Sections~\ref{sec:ch19-scope}--\ref{sec:ch19-spatial-correlators}
used an abstract positive functional on the uncompressed oscillator algebra.
One might now try to interpret its Gaussian contractions as expectation values
in the ground state of the original finite fermionic model.  That identification
is tempting: both constructions have current operators with the same labels,
and the protected low-energy currents satisfy a Schwinger relation on suitable
vectors.  Yet the two constructions have different mathematical carriers and
different hypotheses.  An equality of formal current commutators on a
restricted vector is not an equality of states on the whole operator algebra.

Here is the actual question we can answer.  \emph{Starting from a concrete,
normalized fermionic sea vector, which ordered current products can be
evaluated exactly?  Do these evaluations coincide with the abstract CCR
functional for any choice of the Bogoliubov parameters?}  The first answer
is affirmative in a bounded, explicitly protected regime.  The second is
affirmative for the \emph{free} pair, and cannot be extended by simply
replacing the free pair by an interacting one.  A short calculation below
exhibits the obstruction without taking a large-lattice limit.

To keep the distinction visible, write
\begin{equation}
 \langle X\rangle^{\mathrm F}_{N}:=
 \langle\Omega_N,X\Omega_N\rangle_{\mathcal H_{h}},
 \qquad \omega_P(a):=\omega_0\!\left(\beta_P^{-1}(a)\right).
 \label{eq:ch19-comp-two-evaluators}
\end{equation}
The first expression acts on \emph{finite CAR/Fock operators} $X$;
the second acts on \emph{abstract CCR words} $a$.  Even when the same
letter $\rho$ is used for corresponding currents, these are not, without
an additional map, elements of the same algebra.  A density-matrix state
on the finite carrier is equally possible:
$\langle X\rangle^{\mathrm F}_{\Gamma}
=\operatorname{Tr}_{\mathcal H_h}(\Gamma X)$ for
$\Gamma\geq0$ and $\operatorname{Tr}\Gamma=1$.  The proof below uses the
pure state $\Gamma=|\Omega_N\rangle\langle\Omega_N|$.

\subsection{The finite objects and the protected-current hypotheses}

Fix an integer $h\geq1$ and put $L=2h$.  The finite one-particle band has
$L$ momentum labels per species, and the fermionic Fock space
$\mathcal H_h$ built from this band is finite dimensional.  Let $\sigma$
run through the finite, ordered set of species and fix a charge assignment
$N:\sigma\to\mathbb Z$ admissible for the sea construction in
Chapters~\ref{ch:fermions} and~\ref{ch:vacuum_budget}.  Denote its
normalized sea vector by $\Omega_N$.  In the accepted implementation this
is \texttt{A05.groundKet h Ns}; its normalization is an exact CAR/Fock
statement, not an oscillator assumption:
\begin{equation}
 \langle\Omega_N,\Omega_N\rangle=1,
 \qquad \langle I\rangle_N^{\mathrm F}=1.
 \label{eq:ch19-comp-normalized}
\end{equation}
The term \emph{ground} here refers to the constructed sea or sector ground
for the finite bare setting.  It does not identify this vector as the ground
state of an arbitrary interacting Hamiltonian.

Choose the right species $\nu_R$ in the two-branch record, a positive
mode cutoff $M\in\mathbb N$, and current operators
\begin{equation}
 R_m:=\rho_{m,\nu_R}^{\mathrm F},\qquad
 R_{-m}:=\rho_{-m,\nu_R}^{\mathrm F},\qquad
 R_m^*=R_{-m},\qquad 1\leq m\leq M.
 \label{eq:ch19-comp-current-def}
\end{equation}
The current label is \emph{signed}: $R_m$ creates an excitation of
positive energy grade, whereas $R_{-m}$ lowers it.  It is important not to
replace these signed indices by their absolute values inside an operator
product.  Products act from right to left.

The finite current formulae below assume all of the following:
$h>0$; $N$ admissible; $1\leq m,n\leq M$; and the
\emph{right-branch zero-budget margin condition}
\begin{equation}
 \mathrm{marginAt}(h,M,0,N,\nu_R).
 \label{eq:ch19-comp-margin}
\end{equation}
This is the frozen-margin predicate of the finite implementation, not a
claim about arbitrary lattice parameters.  Its use here is precise: it
ensures that the signed-current commutator theorem applies on $\Omega_N$
for every pair under consideration; in particular its required excursion
inequality
\begin{equation}
 m+n+|N(\nu_R)|\leq h
 \label{eq:ch19-comp-excursion}
\end{equation}
holds for $1\leq m,n\leq M$.  The same construction supplies the
exact lowering property
\begin{equation}
 R_{-m}\Omega_N=0\qquad(1\leq m\leq M).
 \label{eq:ch19-comp-lowering}
\end{equation}
These are \emph{local, vector-dependent facts}.  They do not assert the
operator identity $[R_{-m},R_n]=m\delta_{mn}I$ on all of
$\mathcal H_h$, which would contradict the finite-dimensional trace
obstruction discussed in Section~\ref{sec:ch19-scope}.

Under the margin hypothesis the signed-current theorem gives, \emph{on
the constructed ground vector},
\begin{equation}
 [R_{-m},R_n]\Omega_N=m\delta_{mn}\Omega_N.
 \label{eq:ch19-comp-ground-commutator}
\end{equation}
This is the only commutator relation needed for the protected quadratic
expectations.  No representation of the full uncompressed CCR algebra by
finite matrices is being assumed.

\subsection{Exact ordered contractions in the bare finite ground}

\begin{proposition}[Finite ground-state current contractions]
\label{prop:ch19-comp-ground-contractions}
Under the hypotheses of~\eqref{eq:ch19-comp-margin}, for
$1\leq m,n\leq M$ one has
\begin{align}
 \langle R_{-m}R_n\rangle_N^{\mathrm F}
     &=m\delta_{mn},
 \label{eq:ch19-comp-lr}\\
 \langle R_mR_{-n}\rangle_N^{\mathrm F}
     &=0,
 \label{eq:ch19-comp-rl}\\
 \langle R_mR_n\rangle_N^{\mathrm F}
     &=0,\qquad
 \langle R_{-m}R_{-n}\rangle_N^{\mathrm F}=0.
 \label{eq:ch19-comp-same-sign}
\end{align}
In particular the finite excitation $R_m\Omega_N$ has
$\|R_m\Omega_N\|^2=m$.
\end{proposition}

\begin{proof}
Begin with the first product.  The definition of a commutator,
Equation~\eqref{eq:ch19-comp-ground-commutator}, and the rightmost
lowering relation give
\begin{align*}
 R_{-m}R_n\Omega_N
 &=\bigl([R_{-m},R_n]+R_nR_{-m}\bigr)\Omega_N\\
 &=m\delta_{mn}\Omega_N+R_n(0)
 =m\delta_{mn}\Omega_N.
\end{align*}
Taking the inner product with $\Omega_N$ and using
\eqref{eq:ch19-comp-normalized} proves
\eqref{eq:ch19-comp-lr}.  For the reverse order the lowering mode acts
first, so
$R_mR_{-n}\Omega_N=R_m(0)=0$ and
\eqref{eq:ch19-comp-rl} follows.

Similarly, $R_{-m}R_{-n}\Omega_N=0$.  For the two raising
operators, use the adjoint relation and~\eqref{eq:ch19-comp-lowering}:
\begin{align*}
 \langle\Omega_N,R_mR_n\Omega_N\rangle
 &=\langle R_m^*\Omega_N,R_n\Omega_N\rangle\\
 &=\langle R_{-m}\Omega_N,R_n\Omega_N\rangle=0.
\end{align*}
Finally, adjointness and the already-proved contraction imply
\begin{align*}
 \|R_m\Omega_N\|^2
 &=\langle R_m\Omega_N,R_m\Omega_N\rangle
 =\langle\Omega_N,R_{-m}R_m\Omega_N\rangle=m.
\end{align*}
The argument is finite-dimensional throughout; the nonzero norm is a
certificate for a concrete excited ket in the protected regime.
\end{proof}

Two points are worth separating.  First, the $m$ on the right-hand side
is not an occupation expectation: it is the coefficient of the Schwinger
commutator on the sea vector.  Second, there is no contradiction between
$\langle R_{-m}R_m\rangle=m$ and
$\langle R_mR_{-m}\rangle=0$.  The difference is exactly the
commutator expectation, and its order cannot be changed casually.

\subsection{Finite spatial density and the free-state comparison}

For spatial evaluation let $\Lambda_L=\mathbb Z/L\mathbb Z$ and choose a
primitive $L$th root $\zeta$ (the canonical choice is
$e^{2\pi i/L}$).  With the character
$\chi_m(x)=\zeta^{m x_0}$, where $x_0$ is any integer representative,
introduce the finite right-branch oscillatory density
\begin{equation}
 \rho_R^{\mathrm F}(x):=\frac1L\sum_{m=1}^{M}
 \bigl(\chi_{-m}(x)R_m+\chi_m(x)R_{-m}\bigr).
 \label{eq:ch19-comp-fock-density}
\end{equation}
This operator uses exactly the same selected mode set and Fourier
normalization as the abstract density in
Section~\ref{sec:ch19-spatial-correlators}.  It remains a finite CAR
operator; the equality of its \emph{formula} with a symbolic CCR
expression is not an identification of their representations.
For this character choice the finite density is self-adjoint.

\begin{theorem}[Protected finite-density kernel and free comparison]
\label{thm:ch19-comp-free-kernel}
Assume $h>0$, admissible $N$, the margin condition
\eqref{eq:ch19-comp-margin}, $M\geq0$, and a nonzero root $\zeta$ as
above.  Then for all $x,y\in\Lambda_L$,
\begin{equation}
 \langle \rho_R^{\mathrm F}(x)\rho_R^{\mathrm F}(y)
 \rangle_N^{\mathrm F}
 =\frac{1}{L^2}\sum_{m=1}^{M}m\chi_m(x-y).
 \label{eq:ch19-comp-fock-kernel}
\end{equation}
Moreover, this equals the \emph{abstract} ordered density kernel
evaluated in the free Bogoliubov pair $(c,s)=(1,0)$:
\begin{equation}
 \langle \rho_R^{\mathrm F}(x)\rho_R^{\mathrm F}(y)
 \rangle_N^{\mathrm F}
 =\omega_{\mathrm{free}}\bigl(\rho_R(x)\rho_R(y)\bigr).
 \label{eq:ch19-comp-free-agreement}
\end{equation}
The assertion concerns these quadratic observables and the protected
bare ground, not equality of states on their entire algebras.
\end{theorem}

\begin{proof}
Expand the product in the indicated order:
\begin{align*}
 &\rho_R^{\mathrm F}(x)\rho_R^{\mathrm F}(y)
 =\frac{1}{L^2}\sum_{m,n=1}^{M}\bigl(
 \chi_{-m}(x)\chi_{-n}(y)R_mR_n\\
 &\hspace{15mm}+\chi_{-m}(x)\chi_n(y)R_mR_{-n}
 +\chi_m(x)\chi_{-n}(y)R_{-m}R_n\\
 &\hspace{15mm}+\chi_m(x)\chi_n(y)R_{-m}R_{-n}\bigr).
\end{align*}
Linearity of the finite vector expectation and
Proposition~\ref{prop:ch19-comp-ground-contractions} kill three of
these four products.  The surviving term is the \emph{third}, not the
second:
\begin{align*}
 \langle\rho_R^{\mathrm F}(x)\rho_R^{\mathrm F}(y)
 \rangle_N^{\mathrm F}
 &=\frac{1}{L^2}\sum_{m,n=1}^{M}
 \chi_m(x)\chi_{-n}(y)n\delta_{mn}\\
 &=\frac{1}{L^2}\sum_{m=1}^{M}
 m\chi_m(x)\chi_{-m}(y)\\
 &=\frac{1}{L^2}\sum_{m=1}^{M}m\chi_m(x-y).
\end{align*}
For a merely nonzero complex $\zeta$ the corresponding source identity
is formulated in terms of chosen representatives and integer powers;
primitivity and unit modulus supply the periodic, Hermitian spatial
interpretation here.

The abstract ordered kernel of Section~\ref{sec:ch19-spatial-correlators}
is $L^{-2}\sum_m m\bigl(s^2\chi_{-m}(x-y)
+c^2\chi_m(x-y)\bigr)$.  Substituting $s=0$, $c=1$ leaves
exactly~\eqref{eq:ch19-comp-fock-kernel}.  This proves
\eqref{eq:ch19-comp-free-agreement} by comparing explicitly derived
numbers; no intertwining theorem between the finite CAR algebra and the
abstract CCR algebra has been invoked.
\end{proof}

At coincidence $x=y$ the free kernel gives the exact finite variance
\begin{equation}
 \langle(\rho_R^{\mathrm F}(x))^2\rangle_N^{\mathrm F}
 =\frac{1}{L^2}\sum_{m=1}^{M}m
 =\frac{M(M+1)}{2L^2}.
 \label{eq:ch19-comp-fock-diagonal}
\end{equation}
It is positive for $M\geq1$, and zero for the empty mode set $M=0$.
The evaluation is a genuine nonnegative vector norm since
$\rho_R^{\mathrm F}(x)$ is self-adjoint.  It is \emph{not} an
infinite-volume decay law; both $L$ and $M$ remain finite and explicit.

\subsection{Why the bare ground is not a dressed vacuum}

It remains possible for the \emph{finite} model to have dressed current
combinations, but the existence of those operators does not supply a
vector annihilated by all of them.  Let $\eta\in\{R,L\}$ and
$\bar\eta$ denote the opposite branch.  For a real hyperbolic pair
\begin{equation}
 c,s\in\mathbb R,\qquad c^2-s^2=1,
 \label{eq:ch19-comp-hyperbolic}
\end{equation}
use the finite Bogoliubov operator convention carried by the accepted
implementation:
\begin{align}
 \widetilde A_{\eta m}^{\mathrm F}
 &=c\,\rho_{-m,\eta}^{\mathrm F}
   +s\,\rho_{m,\bar\eta}^{\mathrm F},
 \label{eq:ch19-comp-dressed-ann}\\
 \widetilde C_{\eta m}^{\mathrm F}
 &=c\,\rho_{m,\eta}^{\mathrm F}
   +s\,\rho_{-m,\bar\eta}^{\mathrm F}
 =\bigl(\widetilde A_{\eta m}^{\mathrm F}\bigr)^*.
 \label{eq:ch19-comp-dressed-cre}
\end{align}
Here the signs and the \emph{opposite-branch} mixing are conventions,
not freely adjustable by a generic squeezing formula.  To control the
both-branch matrix elements assume the stronger predicate
$\mathrm{branchMargins}(h,M,0,N,\mathrm{branches})$, so the
protected relations hold for both chosen species.

For a fixed $1\leq m\leq M$ the bare lowering action gives
\begin{equation}
 \widetilde A_{\eta m}^{\mathrm F}\Omega_N
 =s\,\rho_{m,\bar\eta}^{\mathrm F}\Omega_N.
 \label{eq:ch19-comp-ann-action}
\end{equation}
Under the opposite-branch margin its squared norm is $m s^2$.
Thus the obstruction to the dressed vacuum is particularly transparent:
when $s\ne0$, the dressed \emph{annihilator itself creates a nonzero
bare excitation}.  This statement is stronger than merely noticing
that one numerical contraction differs from the abstract result.

\begin{proposition}[Dressed-number witness on the finite bare ground]
\label{prop:ch19-comp-dressed-number}
Under the admissibility, positivity, mode-range, and both-branch margin
hypotheses above,
\begin{align}
 \bigl\langle\widetilde C_{\eta m}^{\mathrm F}
 \widetilde A_{\eta m}^{\mathrm F}\bigr\rangle_N^{\mathrm F}
 &=m s^2,
 \label{eq:ch19-comp-dressed-ca}\\
 \bigl\langle\widetilde A_{\eta m}^{\mathrm F}
 \widetilde C_{\eta m}^{\mathrm F}\bigr\rangle_N^{\mathrm F}
 &=m c^2.
 \label{eq:ch19-comp-dressed-ac}
\end{align}
If $s\ne0$, $\Omega_N$ is not annihilated by the dressed lowering
operator $\widetilde A_{\eta m}^{\mathrm F}$.
\end{proposition}

\begin{proof}
Adjointness, followed by~\eqref{eq:ch19-comp-ann-action}, yields
\begin{align*}
 \bigl\langle\widetilde C_{\eta m}^{\mathrm F}
 \widetilde A_{\eta m}^{\mathrm F}\bigr\rangle_N^{\mathrm F}
 &=\bigl\|\widetilde A_{\eta m}^{\mathrm F}\Omega_N\bigr\|^2\\
 &=s^2\bigl\|\rho_{m,\bar\eta}^{\mathrm F}\Omega_N\bigr\|^2
 =m s^2.
\end{align*}
For the reverse order,
$\widetilde C_{\eta m}^{\mathrm F}\Omega_N
=c\rho_{m,\eta}^{\mathrm F}\Omega_N$, since the opposite-branch
lowering operator kills the bare ground.  Therefore
\begin{align*}
 \bigl\langle\widetilde A_{\eta m}^{\mathrm F}
 \widetilde C_{\eta m}^{\mathrm F}\bigr\rangle_N^{\mathrm F}
 &=\bigl\|\widetilde C_{\eta m}^{\mathrm F}\Omega_N\bigr\|^2
 =c^2\bigl\|\rho_{m,\eta}^{\mathrm F}\Omega_N\bigr\|^2
 =m c^2.
\end{align*}
Since $m>0$, the first value is strictly positive whenever
$s\ne0$, proving that the dressed annihilator cannot kill $\Omega_N$.
Neither step uses a finite global CCR identity.
\end{proof}

The apparent numerical parallel to the abstract contractions requires
care.  In the \emph{abstract interacting} state $\omega_P$, the dressed
annihilators annihilate the state in the algebraic sense, whereas the
bare contractions obey $\omega_P(C_{Rm}A_{Rm})=ms^2$ and
$\omega_P(A_{Rm}C_{Rm})=mc^2$.  Above we obtained the same two
numbers for \emph{dressed} finite operators evaluated in the
\emph{bare} finite ground.  These are different ordered products in
different algebras and different states.  Equality of the displayed
coefficients does not license swapping their roles.  In particular,
$\omega_P(\widetilde C_{Rm}\widetilde A_{Rm})=0$, whereas
\eqref{eq:ch19-comp-dressed-ca} is nonzero for $s\ne0$.

\subsection{A complete finite witness: two modes on sixteen sites}

\begin{example}[Source-verified finite/abstract discrepancy]
\label{ex:ch19-comp-16-sites}
Take $h=8$, so $L=16$, positive-mode cutoff $M=2$, and the explicit
admissible charge/branch data from the CH18 concrete witness.  Their
both-branch margin predicates hold.  The hyperbolic coefficients used
in that witness are
\begin{equation}
 c=\frac54,\qquad s=\frac34,
 \qquad c^2-s^2=\frac{25-9}{16}=1.
 \label{eq:ch19-comp-example-pair}
\end{equation}
For $m=2$, the protected current contraction and excited-vector norm
are
\begin{equation}
 \langle R_{-2}R_2\rangle_N^{\mathrm F}=2,
 \qquad\|R_2\Omega_N\|^2=2,
 \qquad\langle R_2R_{-2}\rangle_N^{\mathrm F}=0.
 \label{eq:ch19-comp-example-bare}
\end{equation}
The dressed-number test gives
\begin{equation}
 \bigl\langle\widetilde C_{R2}^{\mathrm F}
 \widetilde A_{R2}^{\mathrm F}\bigr\rangle_N^{\mathrm F}
 =2\left(\frac34\right)^2=\frac98>0.
 \label{eq:ch19-comp-example-dressed}
\end{equation}
Consequently the \emph{actual finite bare ground} is not an
annihilation vacuum for these nontrivial dressed combinations.

At the same site the finite bare density has
\begin{equation}
 \langle\rho_R^{\mathrm F}(0)\rho_R^{\mathrm F}(0)
 \rangle_N^{\mathrm F}
 =\frac{1+2}{16^2}=\frac3{256}.
 \label{eq:ch19-comp-example-free}
\end{equation}
In contrast, in the \emph{abstract interacting} state the two-mode
ordered kernel is
\begin{align}
 \omega_P\bigl(\rho_R(0)\rho_R(0)\bigr)
 &=\frac{1+2}{16^2}(c^2+s^2)\notag\\
 &=\frac3{256}\frac{34}{16}
 =\frac{51}{2048}.
 \label{eq:ch19-comp-example-abstract}
\end{align}
The two finite numbers are unequal:
$3/256=24/2048\ne51/2048$.  This comparison does \emph{not}
show that no interacting finite ground state exists.  It shows exactly
that the constructed \emph{bare} finite ground and the interacting
abstract quasi-free functional cannot be identified on this chosen
density product.  The first value is certified by the finite CAR
margin calculation; the second by the abstract CCR construction.
\end{example}

Notice that this example uses $c=5/4$ and $s=3/4$.  The distinct
$5/3,4/3$ pair used for the separate compressed-three-dimensional
common-kernel counterexample in Section~\ref{sec:ch19-scope} must not be substituted
here.  These are different mathematical witnesses with different
operators and hypotheses.

\subsection{What an interacting finite comparison would actually require}

The preceding computations settle a useful \emph{free bare-ground
comparison}, but they leave a genuine interacting-state problem.  It
is worth stating the missing pieces because the later order-parameter
chapters will otherwise invite a false shortcut.

\paragraph{A comparison requires an operator map and a specified state.}
Let $\mathfrak A_{\mathrm{CCR}}$ be the algebra used in
Sections~\ref{sec:ch19-scope}--\ref{sec:ch19-spatial-correlators}.
A proposed comparison would have to specify, at least on a chosen
set of words $\mathcal W$, concrete finite-Fock operators
$\Phi_h(a)$ and a finite state $\Gamma_h$.  One could then define
an actual defect, for each \emph{specified} word $a\in\mathcal W$,
\begin{equation}
 \varepsilon_h(a)
 :=\operatorname{Tr}\!\bigl(\Gamma_h\Phi_h(a)\bigr)-\omega_P(a).
 \label{eq:ch19-comp-comparison-defect}
\end{equation}
No theorem in this section establishes that such a map is a unital
$*$-homomorphism on the full CCR algebra; the finite-dimensional trace
obstruction rules out a globally exact nonzero scalar-CCR
representation.  Neither is a bound $\varepsilon_h(a)\to0$ claimed.
For the density products at the free pair we have an explicitly
constructed \emph{wordwise numerical equality}; for arbitrary
interacting pairs no analogous theorem follows from this packet.

\paragraph{Compression introduces a boundary correction.}
Even if one begins with an ambient space $\mathcal H$ and a
finite-rank orthogonal projection $P$ onto an energy-budget subspace,
compression does not commute with operator multiplication.  Put
$Q=I-P$.  The exact algebraic identity on $P\mathcal H$ is
\begin{equation}
 [PXP,PYP]
 =P[X,Y]P-PXQYP+PYQXP.
 \label{eq:ch19-comp-compression-identity}
\end{equation}
It follows by inserting $I=P+Q$ between the factors of
$P[X,Y]P$.  The two last terms measure excursions out of the budget
space and back.  Thus even when an ambient commutator formula is known,
the \emph{compressed} commutator is controlled only after its boundary
remainder has been estimated or shown to vanish on the chosen states.
The protected relation~\eqref{eq:ch19-comp-ground-commutator} is a
local certificate, not a license to discard these terms for every word.

\paragraph{Dressed-vacuum conditions and Hamiltonian ground states
are independent obligations.}
The positive finite operator
$\widetilde C\widetilde A=\widetilde A^*\widetilde A$
is a norm-square test.  The nonzero result
\eqref{eq:ch19-comp-dressed-ca} disproves an annihilation-vacuum
claim for $\Omega_N$; it does not rule out another ground vector of
an interacting finite Hamiltonian $H_{h,\mathrm{int}}$.
To compare that ground to $\omega_P$, one must construct its
normalized vector or density matrix, prove the relevant spectral
statement, and evaluate the same observables with a controlled
projection or comparison mechanism.

\paragraph{Prerequisites for later order parameters.}
Physical CDW, pairing and spinful correlation functions require more
than oscillator two-point contractions: their concrete finite-CAR
operator definitions, charge-sector shifts and Klein-factor signs,
ordering and adjoints, as well as any projection residuals, must be
fixed before assigning a finite expectation.  Formal vertex
exponentials in a power-series algebra are a separate coefficientwise
construction; their identities alone neither define convergent
operators at a numerical parameter nor compare them with finite CAR
observables.  Any claimed continuum power law additionally needs a
specified lattice/cutoff scaling and uniform error estimates.
The CH20 and CH21 developments are not established by the CH19
Phase-B results.

\subsection{Conclusion: exactly what has been compared}

We have constructed the finite expectation
$\langle\cdot\rangle_N^{\mathrm F}$ on a normalized fermionic sea,
used the frozen-margin and lowering hypotheses to prove its ordered
current contractions, and derived an exact finite spatial kernel.
Its agreement with the abstract CCR kernel at $(c,s)=(1,0)$ is a
\emph{quadratic free-state comparison}.  For $s\ne0$ the same bare
ground exhibits a strictly nonzero dressed-number expectation,
providing a direct witness against treating it as the interacting
algebraic vacuum.  The next task is therefore not to repeat the
Gaussian calculation, but to specify and verify a finite interacting
state together with its observable comparison map and boundary
estimates.  The accepted CH19/A08 Phase-B Lean results support the
claims made above; they have not thereby been promoted to the frozen
Core, nor does successful compilation of this prose supply an
independent mathematical review.
